\chapter{The Quine}
% OUTLINE Ch1 "The Quine (self-reproduction earned small)". Laws: both
% inherited (citation + intuition, labeled blocks adds-no-claim); the
% NON-FUNGE seeded (each language's quine is its own). CITATION: Kleene's
% recursion theorem (1938, "On notation for ordinal numbers") as the
% existence guarantee; the name "quine" from Hofstadter 1979 (Godel
% Escher Bach), honoring W.V.O. Quine's use-mention discipline. THE
% KERNEL = the room at its smallest; the "roughly the same size" first
% sighting. Zero device appearances (dose = Ch6 only). Gate: Kodo (the
% Kleene earn).

\begin{intuition}
Ask a machine to print something, and normally you hand it the something
--- a string, kept apart from the machine that prints it. Now ask for the
one output a machine cannot simply be handed: its own complete text,
printed by itself, from nothing but itself. It sounds like asking a
camera to photograph itself in full, lens included, with no mirror. The
feeling to carry: self-reproduction looks paradoxical, right up until it
is built --- and once built it is not magic but a small, exact, and
completely ordinary machine.
\end{intuition}

The machine has a name, given by Douglas Hofstadter in his 1979 book
\emph{G\"odel, Escher, Bach}: a \emph{quine}, after the logician
Willard Van Orman Quine, whose careful discipline about the difference
between using a word and mentioning it is exactly the discipline the
construction turns on. A quine is a program that takes no input and
prints its own source text, character for character. Every general-purpose
programming language admits one, and the reason no reader should find
that surprising is worth walking, because it is this whole volume's engine
in miniature.

The naive attempt fails instructively. Write a program whose instruction
is ``print the following text,'' and paste the program's own text as the
following text --- but now the pasted copy must contain the paste, which
must contain the paste, and the source grows without end before it ever
runs. The trap is exactly Quine's: a text cannot straightforwardly
contain a full quotation of itself, because the quotation is part of the
text it must quote. The escape is the trick every quine shares and the
kernel this chapter is after. Split the program in two. A \emph{body}
that describes, in the ordinary way, how to print some data; and a piece
of \emph{data} that encodes the body. Then run the body twice against the
data: once to print the data as data --- quoted --- and once to print it
as the body it encodes --- executed. The program prints its data, then
prints what its data means, and the two halves together reconstitute the
whole, without any part ever having to contain a full copy of itself. Use
without mention, then mention without use, and the seam closes. That
maneuver --- one description, applied to itself once each way --- is the
gesture the entire book will watch recur.

That such a program must exist, in any system able to describe its own
computations, is not a lucky fact about clever coders; it is a theorem,
and its author is Stephen Kleene, who proved the recursion theorem in
1938. Stated for this chapter's purpose and stripped of its machinery:
any transformation you can describe on programs has a program that
behaves as though handed its own text. Self-reference is not smuggled in;
it is \emph{guaranteed}, for free, by the system's own power to describe
what it does. A language rich enough to talk about its programs cannot
help but contain programs that talk about themselves --- and the quine is
the smallest, plainest witness to that guarantee. The reader should mark
the shape of the promise, because the whole book leans on it: wherever a
system can describe its own operations, the self-referential room exists,
not by permission but by theorem. This chapter has built the door at its
most harmless; the chapters ahead find the same door in houses where its
existence is a shock.

\begin{intuition}
The kernel, felt: hold a short description that says ``write me down, then
carry me out.'' Written down, it is data; carried out, it writes itself
down and carries itself out. The whole self-reproduction lives in that
one dual-purpose sentence, and it is small --- a handful of tokens around
the seam between quoting and doing. That smallness is not incidental. It
is the first sighting of a claim the book will earn in full later: the
room, wherever it appears, is roughly the same size, because it is always
the same little seam.
\end{intuition}

Two honesty notes close the chapter, one guarding a claim already made
and one guarding a claim not yet earned.

The first is the book's own non-funge law, in its gentlest first
instance. Every language has its quine, and they are genuinely
\emph{alike} --- same seam, same use-then-mention maneuver, roughly the
same size. But one language's quine is not another's: it is written in
that language's tokens, runs on that language's rules, and printed into a
foreign interpreter it is not a quine but a syntax error. The likeness is
a family resemblance, exact and real; the identity is false. Even here,
at the smallest room in the book, the rooms do not merge, and the reader
should file the distinction now, because the volume's hardest chapter is
built on it.

The second: the kernel is a mechanism, and this chapter has said nothing
about what it \emph{means}. That a system contains self-describing
programs is a structural fact about the system, cited to Kleene; it is
not, here, a statement about understanding, awareness, or a program
``knowing itself.'' Those readings --- the ones the preface fenced by
name --- add nothing the recursion theorem contains, and this book adds
them nowhere. The quine writes its own text; it does not read it.

What the quine reproduces, it reproduces without judgment --- it never
asks whether its own text is true. The next chapter takes the same
self-application and points it at a question with a truth value, and the
harmless seam becomes the sharpest tool in logic: the diagonal.
